% \begin{document}
\chapter{Operational Semantics}
\dfn{Multi-step evaluation}{
  Binary relation indicated with $ \to^* $ which is the reflexive and transitive closure of the small-step evaluation.
}

Note that the transitive rule is a single step and a multi step, although it can be proved that it allows concatenation of single stepps.

\section{Properties of programs}
\begin{itemize}
\item Determinism: suppose that for any given configuration, if this can step in different ways these are equal
  \item Termination: evaluation terminates. Formally, for the multi step there is a configuration that it can multi step to without other possible steps.
    \item Soundness: if we run a program long enough we should get a final integer, but it's not ture because stores are partial function and if a variable isn't defined we're cooked.
\end{itemize}

\dfn{Free Variables}{
  Given the function $ fvs(e) $ defined inductively as:
  \[ 
  fvs(e) := \begin{cases}
    \emptyset & \text{if } e \text{ is a variable} \\
    fvs(e_1) \cup \{ x \} & \text{if } e = (\lambda x. e_1) \\
    fvs(e_1) \cup fvs(e_2) & \text{if } e = (e_1 e_2)
  \end{cases}
  \]
}


Note that in $ x \coloneq x + 1; y \coloneq x $ we still need to add $ x $ to the free variables because it's in the rhs of the first assignment. 

\dfn{Well-Formedness}{
  For any environment $ \sigma $ and expression $ e $, if all free variables in $ e $ are in $ \sigma $, then $ e $ is well-formed under $ \sigma $.
}

Two properties that imply soundness:

\dfn{Progress}{
  Every well-formed configuration is either an \texttt{Int} or can step to another configuration.
}

\dfn{Preservation}{
  If a configuration is well-formed, any following configuration it stepps to will also be well-formed.
}

with these lenmas and termination we can show that a well-formaed configuration is sound.

\subsection{Inductive Sets}
A bit of maths before the proof

\dfn{Inductive Set}{
A set $S$ is \textbf{inductive} if:
\begin{enumerate}
    \item $\emptyset \in S$, and
    \item for every set $x$, if $x \in S$, then $x \cup \{x\} \in S$.
\end{enumerate}
}

We have used inductive setts with the inductive rules (remember that $ \to $ is a relation, so a set of tuples) and BNF grammar, free-variables and multistep. Another example is the natural numbers using the successor definition.

\subsection{Inductive Proofs}
If we want to prove something about an inductive set (that can be infinite by closing under the inductive rules) we can use the induction principle.

The base cases are all the axioms.

Sometimes the properties are defined over multiple inductive setts, so we need to chose one. 
% \end{document}
